\documentclass[11pt]{article}
\usepackage[utf8]{inputenc}
\usepackage[T1]{fontenc}
\usepackage[margin=1in]{geometry}
\usepackage{hyperref}
\usepackage{enumitem}
\title{FLEX Week | Or Digital Tools/Purpose}
\author{Piter Garcia}
\date{March 20, 2026}
\begin{document}
\maketitle
\section*{Quick Scan}
\begin{itemize}[leftmargin=*]
\item Grade: Grade 4
\item Workbook sessions: Session 14
\item Class blocks: Day 1 1:15-2:10 Baris, Day 2 1:15-2:10 Fowler, Day 3 1:15-2:10 Schrank, Day 4 1:15-2:10 Autore
\item Reference file: flex-week-or-digital-tools-purpose.bib
\end{itemize}
\section*{School Planning Goals and Inclusion}
\begin{itemize}[leftmargin=*]
\item What is expected: - I can select the tool that best helps me communicate and collaborate with others based on my task/purpose. | - I can note common features among different digital technology tools. | - I can reason about when to keep information private and when to share information.
\item What we achieved: No direct transcript match is attached yet. Treat this lesson as workbook-backed and enrich it when a matching classroom transcript is available.
\item What needs improvement: stronger proof, cleaner assessment notes, and tighter lesson artifacts.
\item How to improve: add transcript proof, one saved artifact, and clearer success criteria.
\item Inclusion focus: use visual modeling, chunked directions, predictable routines, and flexible participation roles.
\end{itemize}
\section*{Official References}
\begin{itemize}[leftmargin=*]
\item \url{https://csteachers.org/k12standards/interactive/}

\end{itemize}
\end{document}
